\documentclass[11pt]{article}

\usepackage[margin=1in]{geometry}
\usepackage{amsmath, amssymb}
\usepackage{graphicx}
\usepackage{booktabs}
\usepackage{hyperref}
\usepackage{listings}
\usepackage{xcolor}
\usepackage{algorithm}
\usepackage{algpseudocode}
\usepackage{caption}

\hypersetup{colorlinks=true, linkcolor=blue, citecolor=blue, urlcolor=blue}

\lstset{
    basicstyle=\ttfamily\small,
    backgroundcolor=\color{gray!10},
    frame=single,
    breaklines=true,
    columns=fullflexible,
    showstringspaces=false
}

\title{\textbf{Deep Bluffing: A Heads-Up Limit Hold'em Poker Agent} \\
\large ICG Summer Project 2026--27 --- Capstone Report}
\author{Aryaa \\ Roll Number: 250210 }
\date{\today}

\begin{document}

\maketitle

\begin{abstract}
This report describes the design and implementation of an autonomous agent for
Heads-Up Limit Texas Hold'em (HULHE), built entirely from scratch in Python
without reliance on external poker libraries. The system consists of four
layered components: (i) an exact combinatorial hand evaluator built on
\texttt{itertools.combinations} and \texttt{collections.Counter}, (ii) a fast
preflop hand-strength estimator based on the Chen Formula, (iii) a Monte
Carlo rollout estimator for postflop win probability, and (iv) a
Counterfactual Regret Minimization (CFR) strategy trained offline over a
card-abstracted version of the game. Decisions are made primarily by the
trained CFR strategy, with the Monte Carlo / expected-value heuristic serving
as a provably safe fallback for any game state not covered during training.
Local simulation against a random baseline opponent shows a positive
win rate of roughly $+0.8$ big blinds per hand.
\end{abstract}

\tableofcontents

\section{Introduction}

Heads-Up Limit Hold'em is a well-studied benchmark in imperfect-information
game theory: unlike perfect-information games such as Tic-Tac-Toe or the
CartPole/FrozenLake environments used in earlier assignments of this course,
a poker agent must act under genuine uncertainty about its opponent's private
information. The goal of this project was to build both the game engine
primitives (card representation, hand evaluation) and a competitive decision
policy, without using pre-built poker libraries such as \texttt{treys} or
\texttt{RLCard}.

The final agent is organized into independent, testable modules:

\begin{itemize}
    \item \texttt{card.py} --- a \texttt{Card} class with rank/suit parsing and
    comparison dunder methods.
    \item \texttt{deck.py} --- a standard 52-card deck with dealing utilities.
    \item \texttt{hand\_evaluator.py} --- exact 5-card and best-of-7-card hand
    ranking.
    \item \texttt{agent.py} --- the required \texttt{BasePokerBot} /
    \texttt{CustomPokerBot} entry point, containing the full decision logic.
    \item \texttt{cfr\_trainer.py} --- an offline script that trains a CFR
    strategy over an abstracted version of the game and serializes it to
    \texttt{cfr\_strategy.json}.
\end{itemize}

\section{Hand Evaluation}

Rather than hand-writing comparison logic between every pair of hand
categories (flush vs. straight, full house vs. flush, etc.), each 5-card hand
is scored as a Python tuple
\[
    \text{score} = (\text{category}, \text{tiebreak}_1, \text{tiebreak}_2, \dots)
\]
where \emph{category} $\in \{0, \dots, 8\}$ ranges from High Card to Straight
Flush. Python's native tuple comparison then handles all comparisons
correctly and automatically: the category dominates the comparison, and ties
within a category fall through to the kicker values. For example, a full
house is encoded as $(6, r_{\text{trips}}, r_{\text{pair}})$, so any full
house tuple is guaranteed to exceed any flush tuple $(5, \dots)$ regardless of
the specific ranks involved.

Given 7 cards (2 hole cards + 5 community cards), the best possible 5-card
hand is found via
\[
    \texttt{best\_hand}(C) = \max_{S \in \binom{C}{5}} \texttt{evaluate\_5}(S),
\]
implemented with \texttt{itertools.combinations(cards, 5)} over all 21
combinations, and \texttt{collections.Counter} to detect pairs, trips,
quads, and full houses from rank frequency counts.

\paragraph{Correctness.} The evaluator was verified with a dedicated test
suite (\texttt{test\_hand\_evaluator.py}) covering all 9 hand categories, the
wheel straight (A-2-3-4-5, where the Ace plays low), kicker-based tie
resolution, and exact split-pot detection --- all of which pass.

\section{Preflop Strategy: The Chen Formula}

A Monte Carlo rollout is a poor tool for preflop decisions: with only 2 of 7
relevant cards known, a rollout of the remaining 5 board cards plus the
opponent's 2 hole cards is both slow and has high variance for a fixed
simulation budget. Instead, preflop strength is estimated using the
\textbf{Chen Formula}, a well-known closed-form heuristic used informally by
professional players to rank the 169 canonical starting hands. Given the
high card $h$ and low card $l$ of a starting hand:

\begin{enumerate}
    \item \textbf{Base points} from the high card: Ace $=10$, King $=8$, Queen
    $=7$, Jack $=6$; otherwise $\text{rank}(h)/2$.
    \item \textbf{Pair bonus}: if $h = l$ (a pocket pair), points are doubled,
    floored at a minimum of 5.
    \item \textbf{Suited bonus}: $+2$ if $h$ and $l$ share a suit.
    \item \textbf{Gap penalty}: for non-pairs, a penalty is subtracted based on
    the rank gap $g = \text{rank}(h) - \text{rank}(l) - 1$: $-1$ for $g=1$,
    $-2$ for $g=2$, $-4$ for $g=3$, $-5$ for $g \geq 4$.
    \item \textbf{Connector bonus}: $+1$ if $g \leq 1$ and $h < $ Queen (extra
    straight potential for low/medium connectors).
\end{enumerate}

This yields a score roughly in $[0, 20]$, which is then linearly mapped onto
an approximate heads-up win probability:
\[
    P(\text{win}) = 0.30 + 0.55 \times \frac{\min(\text{score}, 20)}{20},
\]
calibrated so that pocket Aces ($\text{score}=20$) map to $P \approx 0.85$
and the weakest hands map to the empirically observed heads-up floor of
$P \approx 0.30$ (in heads-up play, even the worst starting hand retains
meaningful equity against a random opponent hand, unlike in full-ring play).
Table \ref{tab:chen} shows this ranking is directionally correct against
poker intuition.

\begin{table}[h]
\centering
\begin{tabular}{lccc}
\toprule
Hand & Description & Chen Score & Estimated $P(\text{win})$ \\
\midrule
AA & Pocket Aces & 20.0 & 0.85 \\
KK & Pocket Kings & 16.0 & 0.74 \\
AKs & Ace-King suited & 12.0 & 0.63 \\
AKo & Ace-King offsuit & 10.0 & 0.57 \\
98s & Suited connector & 7.5 & 0.51 \\
22 & Lowest pocket pair & 5.0 & 0.44 \\
72o & Worst starting hand & 0.0 & 0.30 \\
\bottomrule
\end{tabular}
\caption{Chen Formula scores and mapped win probabilities for representative hands.}
\label{tab:chen}
\end{table}

This approach reduces preflop decision latency from tens of milliseconds
(Monte Carlo) to sub-millisecond, which matters when the tournament arena
runs 10{,}000 hands with multiple decision points each.

\section{Postflop Strategy: Monte Carlo Rollouts}

Once community cards are visible, there is genuine information to exploit,
so postflop win probability is estimated directly via Monte Carlo rollout.
For $N=100$ trials, the estimator samples a random completion of the board
and a random opponent hand from the remaining unseen cards, evaluates both
hands with \texttt{best\_hand}, and accumulates:
\[
    \hat{P}(\text{win}) = \frac{1}{N}\sum_{i=1}^{N}
    \begin{cases}
        1 & \text{if my hand beats the sampled opponent hand} \\
        0.5 & \text{if tied} \\
        0 & \text{otherwise}
    \end{cases}
\]
This is an unbiased Monte Carlo estimator of the true win probability
conditional on the known cards; its standard error scales as
$O(1/\sqrt{N})$, which is why $N=100$ was chosen as a practical trade-off
between estimate stability and the $\sim$10{,}000-hand runtime budget.

\section{Decision Rule: Pot Odds and Expected Value}

Given a win probability estimate $\hat{P}$, the baseline decision rule
compares it against the pot odds required to profitably continue:
\[
    P_{\text{required}} = \frac{\text{amount\_to\_call}}{\text{pot\_size} + \text{amount\_to\_call}}.
\]
If $\hat{P} > P_{\text{required}}$, calling has positive expected value:
\[
    \mathbb{E}[\text{EV}_{\text{call}}] = \hat{P} \cdot (\text{pot\_size} + \text{amount\_to\_call}) - (1-\hat{P}) \cdot \text{amount\_to\_call} > 0.
\]
Hands with $\hat{P} > 0.75$ raise for value; hands below $P_{\text{required}}$
(with a small $0.05$ tolerance buffer to absorb Monte Carlo noise) fold.
When \texttt{amount\_to\_call} $= 0$ (a free check), the agent never folds,
since checking has zero downside.

\subsection{Controlled Bluffing}

A purely EV-based policy is exploitable: an opponent that observes many
hands can learn that this agent only continues with genuine equity, and can
profitably bet into it with any two cards on later streets. To counter this,
whenever the EV heuristic would otherwise fold to a bet that costs less than
20\% of the current stack, the agent raises instead with probability $0.08$
(and folds otherwise). This bounded, low-frequency bluff limits the maximum
damage a bad bluff can do while still making the agent's behavior
non-deterministic and harder to exploit --- consistent with the fact that
optimal poker strategies are inherently mixed strategies rather than pure
threshold rules.

\section{Counterfactual Regret Minimization}

\subsection{Motivation and Abstraction}

CFR (Zinkevich et al., 2007) is an iterative self-play algorithm that
converges to a Nash equilibrium strategy in two-player zero-sum
imperfect-information games. Applied to the exact card space of HULHE, the
state space is intractably large (billions of information sets). Real
production poker agents such as Libratus and Pluribus address this with
\emph{card abstraction}: grouping strategically similar hands into a small
number of buckets and solving CFR exactly over the abstracted game. This
project adopts the same principle at a scale appropriate for a class
project.

\paragraph{Bucketing.} Each player's hand strength at any decision point is
mapped to one of 5 buckets via
\[
    \text{bucket} = \min(\lfloor 5 \cdot \hat{P}(\text{win}) \rfloor, 4) \in \{0, 1, 2, 3, 4\},
\]
using the same win-probability estimator described above (Chen Formula
preflop, Monte Carlo postflop). Bucket 0 represents the weakest hands and
bucket 4 the strongest.

\paragraph{Street-to-street evolution.} During offline training, a player's
bucket is not tied to literal cards; instead, it evolves as a bounded random
walk between streets,
\[
    b_{t+1} = \text{clip}(b_t + \delta, 0, 4), \quad
    \delta = \begin{cases} -1 & \text{w.p. } 0.2 \\ 0 & \text{w.p. } 0.6 \\ +1 & \text{w.p. } 0.2 \end{cases}
\]
reflecting the intuition that a hand's relative strength drifts somewhat as
community cards are revealed, without the computational cost of tracking
literal card combinations through the training loop.

\paragraph{Showdown resolution.} Since exact cards are not tracked during
training, the probability that bucket $b_A$ beats bucket $b_B$ at showdown is
approximated as a monotonic function of the bucket gap $\Delta = b_A - b_B$
(Table \ref{tab:showdown}), in the spirit of an Elo-style win-probability
curve: a larger gap implies a higher win probability, and equal buckets
imply a fair (50/50 in expectation) contest.

\begin{table}[h]
\centering
\begin{tabular}{cccccccccc}
\toprule
$\Delta$ & $-4$ & $-3$ & $-2$ & $-1$ & $0$ & $+1$ & $+2$ & $+3$ & $+4$ \\
\midrule
$P(A \text{ wins})$ & 0.05 & 0.15 & 0.25 & 0.35 & 0.50 & 0.65 & 0.75 & 0.85 & 0.95 \\
\bottomrule
\end{tabular}
\caption{Approximate showdown win probability as a function of bucket gap.}
\label{tab:showdown}
\end{table}

\paragraph{Betting abstraction.} Only the \emph{cards} are abstracted --- the
betting structure is preserved exactly as specified: fixed bet sizes of 2
chips (preflop/flop) and 4 chips (turn/river), a cap of 4 total bets per
street, and the specified heads-up turn order.

\subsection{Algorithm: External-Sampling Monte Carlo CFR}

Vanilla CFR requires a full traversal of the game tree on every iteration,
which is unnecessary here since the abstracted game is small. Instead we use
\textbf{External-Sampling Monte Carlo CFR} (Lanctot et al., 2009): on each
iteration, the traversing player explores \emph{all} of their legal actions
(needed to compute counterfactual regret), while the opponent's actions and
chance events (bucket transitions) are each sampled once according to the
current strategy / transition distribution. This is far cheaper per
iteration than vanilla CFR while retaining the same convergence guarantee in
expectation.

At each information set (keyed by \texttt{(bucket, street, bets\_count,
facing\_bet)}), the strategy is obtained via \textbf{regret matching}:
\[
    \sigma(a) = \frac{\max(R(a), 0)}{\sum_{a'} \max(R(a'), 0)}
    \quad \text{(or uniform if all regrets are non-positive)},
\]
where $R(a)$ is the cumulative counterfactual regret for action $a$. After
training, the \emph{average} strategy (not the final iteration's strategy)
is extracted and used at inference time, since the time-averaged strategy is
the one with equilibrium convergence guarantees in CFR.

Training was run for 300{,}000 iterations, converging to 100 distinct
information sets (the full size of the abstracted infoset space) in
approximately 154 seconds on a single CPU core. Table
\ref{tab:strategy-sample} shows a representative slice of the learned
strategy.

\begin{table}[h]
\centering
\begin{tabular}{lccc}
\toprule
Bucket & Fold & Call & Raise \\
\midrule
0 (weakest) & 1.00 & 0.00 & 0.00 \\
1 & 0.00 & 1.00 & 0.00 \\
2 & 0.00 & 1.00 & 0.00 \\
3 & 0.00 & 1.00 & 0.00 \\
4 (strongest) & 0.00 & 0.00 & 1.00 \\
\bottomrule
\end{tabular}
\caption{Learned river strategy facing a bet, by bucket. The strategy is
strongly monotonic in hand strength, as expected: the weakest bucket always
folds and the strongest always raises for value.}
\label{tab:strategy-sample}
\end{table}

Interestingly, the learned \emph{preflop} strategy facing the blind
occasionally favors calling (``limping'') over raising even for
moderate-to-strong buckets. While counter-intuitive at first glance, this is
a legitimate emergent property of the abstracted game: with only 5 discrete
buckets and a bet cap of 4, immediately raising forecloses future betting
rounds' information value, and the equilibrium strategy in this simplified
game can favor delaying aggression. This is a known phenomenon in CFR
research --- equilibrium strategies are frequently non-obvious and do not
always match single-street heuristic intuition, which is part of what makes
game-theoretic solutions valuable over hand-coded rules.

\subsection{Inference-Time Integration}

At decision time, the current game state is mapped to an information-set key
and looked up in the trained strategy table. If found, the agent
\textbf{samples} an action proportionally to the trained probabilities
(rather than taking the arg-max), since playing the exact mixed strategy is
the theoretically correct way to deploy an average CFR strategy. If the
infoset is not present in the table --- or the strategy file fails to load
for any reason --- the agent falls back to the EV heuristic of Section 5,
guaranteeing the agent never crashes and always returns a legal action.

\paragraph{A practical limitation.} The \texttt{get\_action()} interface
mandated by the assignment does not expose the betting history of the
current street (i.e., how many bets/raises have already occurred). Since the
bot instance persists across all decisions within a hand and heads-up turns
strictly alternate, this is reconstructed internally: if the agent is asked
to act again on the same street without the board changing, this can only
mean the opponent raised (a call or fold would have ended the street or
hand). Combined with tracking the agent's own chosen actions, this yields an
accurate --- though not perfectly exact for the very first decision after an
opponent's opening action on a street where the agent acts second ---
reconstruction of the betting count used for the infoset lookup. This
approximation is a deliberate, documented trade-off rather than an oversight.

\section{Results}

A standalone local game engine (\texttt{simulate.py}, not part of the
submission) implementing the exact assignment rules was used to evaluate the
agent against a random baseline opponent (which calls most actions, folding
or raising with low fixed probability). Over 2{,}000 simulated hands:

\begin{itemize}
    \item Net result: approximately $+0.8$ big blinds per hand on average.
    \item The agent reliably folds weak hands preflop and postflop, raises
    strong made hands (e.g. always raising a river nut flush), and never
    folds when facing a free check.
\end{itemize}

While a random baseline is a weak benchmark, it confirms the basic
soundness of the decision pipeline end-to-end (evaluator $\to$ probability
estimate $\to$ CFR/EV decision) before deployment into the tournament arena
against peer-submitted agents.

\section{Limitations and Future Work}

\begin{itemize}
    \item \textbf{Bucket granularity.} Five buckets is a coarse abstraction;
    a natural extension is increasing bucket count, or using a continuous
    equity value with function approximation (e.g. a small neural network)
    instead of discrete buckets.
    \item \textbf{Independent bucket evolution.} The random-walk bucket
    transition model treats each player's hand strength as evolving
    independently, when in reality both players' equities are correlated
    through the shared community cards. A joint transition model, estimated
    empirically via sampling, would be more faithful.
    \item \textbf{Betting history reconstruction.} As discussed in Section
    6.3, the bet-count reconstruction is exact in most cases but can be
    approximate on a player's first decision of a street when they act
    second. Exposing betting history explicitly in the interface (or
    tracking a hand ID) would remove this limitation entirely.
    \item \textbf{Opponent modeling.} The current agent does not adapt to a
    specific opponent's observed tendencies within a match. An extension
    using online opponent exploitation (e.g. adjusting away from equilibrium
    once a leak is detected) could improve performance against weak
    opponents in the tournament, at the cost of exploitability against
    strong ones.
\end{itemize}

\section{Conclusion}

This project implements a complete, from-scratch HULHE game engine and
decision agent, layering an exact hand evaluator, a closed-form preflop
heuristic (Chen Formula), Monte Carlo postflop equity estimation, and an
abstracted CFR-trained strategy with a provably safe EV-heuristic fallback.
Each layer is independently unit-tested, and the complete pipeline was
validated via a from-scratch local game-engine simulation prior to
submission. The CFR component, while trained over a necessarily simplified
abstraction, reproduces qualitatively correct poker behavior (monotonic
fold/call/raise thresholds by hand strength) and illustrates the core
game-theoretic principles --- regret minimization and equilibrium mixed
strategies --- that underlie state-of-the-art poker agents.



\end{document}
